\textbf{Variant 1.} A grasshopper jumps along the real axis. He starts at point $0$ and makes $2009$ jumps to the right with lengths $1, 2, \ldots, 2009$ in an arbitrary order. Let $M$ be a set of $2008$ positive integers less than $1005 \cdot 2009$. Prove that the grasshopper can arrange his jumps in such a way that he never lands on a point from $M$.

\textbf{Variant 2.} Let $n$ be a nonnegative integer. A grasshopper jumps along the real axis. He starts at point $0$ and makes $n+1$ jumps to the right with pairwise different positive integral lengths $a_1, a_2, \ldots, a_{n+1}$ in an arbitrary order. Let $M$ be a set of $n$ positive integers in the interval $(0, s)$, where $s = a_1 + a_2 + \cdots + a_{n+1}$. Prove that the grasshopper can arrange his jumps in such a way that he never lands on a point from $M$.
